\documentclass[a4paper,11pt]{article}

\newcommand{\TPnumero}{Lab 3 -- ACF / PACF}
% =========================================================
% Tutorial / lab preamble — Time Series Analysis module
% article class + fancyhdr + tcolorbox + listings (English)
% Author : Aymen Ben Brik (aymen.benbrik@esprit.tn)
% =========================================================

\usepackage[utf8]{inputenc}
\usepackage[T1]{fontenc}
\usepackage[english]{babel}
\usepackage{lmodern}
\usepackage{amsmath, amssymb, amsthm, mathtools, bm}
\usepackage{graphicx}
\usepackage{booktabs}
\usepackage{array}
\usepackage{multirow}
\usepackage{colortbl}
\usepackage{xcolor}
\usepackage{tikz}
\usetikzlibrary{positioning, arrows.meta, calc, shapes, fit, backgrounds, matrix}
\usepackage{listings}
\usepackage[most]{tcolorbox}
\usepackage{geometry}
\geometry{a4paper, margin=2cm, headheight=15pt}
\usepackage{fancyhdr}
\usepackage[hidelinks]{hyperref}
\usepackage{enumitem}

% --- Colours ---
\definecolor{espritBlue}{RGB}{30, 60, 110}
\definecolor{espritAccent}{RGB}{200, 80, 30}
\definecolor{boxEx}{RGB}{30, 60, 110}
\definecolor{boxQ}{RGB}{180, 120, 0}
\definecolor{boxC}{RGB}{30, 100, 60}
\definecolor{boxAtt}{RGB}{200, 80, 30}

% --- Listings ---
\definecolor{codebg}{RGB}{248,248,248}
\definecolor{codeframe}{RGB}{220,220,220}
\definecolor{keyword}{RGB}{0,82,155}
\definecolor{comment}{RGB}{88,110,117}
\definecolor{string}{RGB}{133,153,0}

\lstdefinestyle{pythonstyle}{
  language=Python,
  basicstyle=\ttfamily\small,
  keywordstyle=\color{keyword}\bfseries,
  commentstyle=\color{comment}\itshape,
  stringstyle=\color{string},
  backgroundcolor=\color{codebg},
  frame=single,
  rulecolor=\color{codeframe},
  frameround=tttt,
  breaklines=true,
  showstringspaces=false,
  tabsize=2
}
\lstdefinestyle{Rstyle}{
  language=R,
  basicstyle=\ttfamily\small,
  keywordstyle=\color{keyword}\bfseries,
  commentstyle=\color{comment}\itshape,
  stringstyle=\color{string},
  backgroundcolor=\color{codebg},
  frame=single,
  rulecolor=\color{codeframe},
  frameround=tttt,
  breaklines=true,
  showstringspaces=false,
  tabsize=2
}
\lstset{style=pythonstyle}

% --- Common macros (configurable path) ---
\providecommand{\macrosPath}{../../preamble/macros.tex}
\input{\macrosPath}

% --- Exercise counter ---
\newcounter{excounter}
\renewcommand{\theexcounter}{\arabic{excounter}}

\newtcolorbox[use counter=excounter]{exercise}[1][]{
  enhanced, breakable,
  colback=boxEx!5, colframe=boxEx, fonttitle=\bfseries,
  title={Exercise~\theexcounter\ifstrempty{#1}{}{ -- #1}},
  arc=2pt, boxrule=0.6pt
}
\newtcolorbox{question}[1][]{
  enhanced, breakable,
  colback=boxQ!5, colframe=boxQ, fonttitle=\bfseries,
  title={Question\ifstrempty{#1}{}{ -- #1}},
  arc=2pt, boxrule=0.5pt
}
\newtcolorbox{solution}[1][]{
  enhanced, breakable,
  colback=boxC!5, colframe=boxC, fonttitle=\bfseries,
  title={Solution\ifstrempty{#1}{}{ -- #1}},
  arc=2pt, boxrule=0.5pt
}
\newtcolorbox{warning}[1][]{
  enhanced, breakable,
  colback=boxAtt!5, colframe=boxAtt, fonttitle=\bfseries,
  title={Warning\ifstrempty{#1}{}{ -- #1}},
  arc=2pt, boxrule=0.5pt
}

% --- Headers / footers ---
\pagestyle{fancy}
\fancyhf{}
\fancyhead[L]{\textsc{\moduleNom} -- \auteurNom}
\fancyhead[R]{\TPnumero}
\fancyfoot[C]{\thepage}
\fancyfoot[L]{\auteurInstitution}
\fancyfoot[R]{\auteurEmail}
\renewcommand{\headrulewidth}{0.4pt}
\renewcommand{\footrulewidth}{0.2pt}

% --- Placeholder ---
\providecommand{\TPnumero}{Lab}


\title{Lab 3: ACF, PACF, and white-noise tests\\
\large Sample autocorrelations on simulated and real series}
\author{\auteurNom\\\auteurEmail\\\auteurInstitution}
\date{\anneeUniversitaire}

\begin{document}
\maketitle

\section*{Goals}
\begin{itemize}
  \item Compute and visualise the sample ACF and PACF.
  \item Verify Bartlett's bands on a simulated white noise.
  \item Check the theoretical ACF/PACF signatures of AR(1) and MA(1)
        against simulations.
  \item Apply ACF/PACF to the Atlanta series and run the Ljung--Box
        test to assess white-noise residuals.
\end{itemize}

\section*{Setup}
\begin{itemize}
  \item Python 3 with \texttt{numpy}, \texttt{statsmodels},
        \texttt{matplotlib}.
  \item \texttt{numpy.random.seed(42)}.
  \item Dataset: \texttt{data/AvTempAtlanta.txt} (1656 monthly obs).
\end{itemize}

% =========================================================
\begin{exercise}[White noise simulation and ACF]
\textbf{(1)} Simulate a Gaussian white noise of length $T = 500$:
$\varepsilon_t \iid \mathcal{N}(0, 1)$.

\textbf{(2)} Plot $\varepsilon_t$ vs.\ $t$. Plot the sample ACF up to
lag $30$ with the $\pm 1.96 / \sqrt{T}$ Bartlett bands.

\textbf{(3)} Count the proportion of lags $h \in \{1, \dots, 30\}$
where $\abs{\widehat\rho(h)} > 1.96 / \sqrt{T}$. Under $H_0$ this should
be close to $5\%$. Comment.
\end{exercise}

% =========================================================
\begin{exercise}[AR(1) simulation and theoretical ACF]
\textbf{(1)} For $\phi \in \{0.3, 0.7, -0.7, 0.95\}$, simulate
$X_t = \phi X_{t-1} + \varepsilon_t$ of length $T = 500$ with
$\varepsilon_t \iid \mathcal{N}(0, 1)$, $X_0 = 0$.

\textbf{(2)} For each $\phi$, plot the empirical ACF up to lag 20 and
overlay $\rho_{\text{theo}}(h) = \phi^h$ as red points.

\textbf{(3)} Plot the empirical PACF for $\phi = 0.7$. Is it cut off
at lag 1, as the theory predicts?

\textbf{(4)} For $\phi = 0.95$, the series is \emph{nearly} non-stationary.
What does the empirical ACF look like? How does it differ from the
theoretical $0.95^h$ on a finite sample of size 500?
\end{exercise}

% =========================================================
\begin{exercise}[MA(1) simulation]
\textbf{(1)} For $\theta \in \{0.5, -0.5, 0.9\}$, simulate
$X_t = \varepsilon_t + \theta \varepsilon_{t-1}$ of length $T = 500$.

\textbf{(2)} Plot the empirical ACF and verify that:
\begin{itemize}
  \item $\widehat\rho(1) \approx \theta / (1 + \theta^2)$;
  \item $\widehat\rho(h) \approx 0$ for $h \geq 2$.
\end{itemize}

\textbf{(3)} Plot the empirical PACF: is it decaying as expected?
\end{exercise}

% =========================================================
\begin{exercise}[ACF / PACF on Atlanta temperatures]
Load \texttt{data/AvTempAtlanta.txt} and stack to a monthly series
($n = 1656$).

\textbf{(1)} Plot the raw ACF up to lag 60. What do you see?
Justify the strong spikes at lags 12, 24, 36.

\textbf{(2)} Decompose the series (Chapter 2, additive workflow with
$2 \times 12$ MA + seasonal averaging) and compute the residual
$\widehat\varepsilon_t$.

\textbf{(3)} Plot the ACF and PACF of $\widehat\varepsilon_t$. Are
the spikes at lags 12, 24, 36 still present? What does this say about
the quality of the decomposition?
\end{exercise}

% =========================================================
\begin{exercise}[Ljung--Box test]
\textbf{(1)} Use \texttt{statsmodels.stats.diagnostic.acorr\_ljungbox}
to test, on the residual $\widehat\varepsilon_t$ of Exercise 4:
$H_0$: residuals are white noise, vs.\ $H_1$: residuals have
autocorrelation, at lags $m \in \{12, 24, 36\}$ and significance level
$\alpha = 5\%$.

\textbf{(2)} Report the test statistic, $p$-value, and conclusion at
each lag. What does this tell you about the residual?

\textbf{(3) [bonus]} Apply the same test to the simulated white noise
of Exercise 1 — you should \emph{not} reject $H_0$.
\end{exercise}

\bigskip
\begin{warning}
The functions \texttt{statsmodels.graphics.tsaplots.plot\_acf} and
\texttt{plot\_pacf} draw the Bartlett bands automatically.
\end{warning}

\end{document}
